% GDZ PPN236018647, native495 middle; HKUST496; original diary2v.
% Entries corresponding to printed34-38; adjacent Sept14 retained in its manuscript row.
\clearpage\SourceLeaf{495}\NoSourceHead
\begingroup\fontsize{10}{15}\selectfont\itshape
\newcommand{\DiaryStrike}[1]{\tikz[baseline=(t.base)]{\node[inner sep=0pt,outer sep=0pt] (t) {\strut #1};\draw[line width=.55pt] ([yshift=.15ex]t.west)--([yshift=.15ex]t.east);}}
\noindent Methodus facilis inveniendi aeq: in \(y\) ex\hfill Sept. 14\\
\hspace*{2em}aeq. in \(x\) si ponatur \(x^n+ax^{n-1}+bx^{n-2}\ldotp\,\ldotp\,\ldotp=y\)

\medskip
\noindent fractiones quarum denominator continet quantitates irrationa-\\
\hspace*{2em}les (quomodocunque?) in alias transmutare\\
\hspace*{2em}ab hoc incommodo liberatas.\hfill Sept. 16

\medskip
\noindent\underline{Coefficientes aeq: auxiliariae eliminationi inservientis}\\
\hspace*{2em}\underline{ex radicibus aeq: datae determinati}\hfill eod.

\medskip
\noindent Nova methodus qua resolutionem aequationum universalem\\
\hspace*{2em}investigare forsitanque invenire licebit\hfill Sept. 17.\\
\hspace*{2em}Scil. transm: aeq: in aliam cuius radices
\[
\alpha\varrho'+\beta\varrho''+\gamma\varrho'''+\ldotp\,\ldotp\,\ldotp
\quad\text{ubi}\quad{}^n\!\surd1=\alpha,\beta\ \text{pp.}\quad\&\quad n\ \text{numerus}
\]
\hspace*{2em}aequationis gradum denotans

\medskip
\noindent\underline{In mentem mihi venit radices aeq. \(x^n-1\) ex aeq., communes}\\
\hspace*{2em}\underline{radices habentibus elicere ut adeo plerumque tantum}\\
\hspace*{2em}\underline{aequationes coefficientibus rationalibus gaudent\DiaryStrike{ibus}es}\\
\hspace*{2em}\underline{resolvi oporteat}\hfill Sept. 29 Bruns.

\medskip
\endgroup
